\section{Correctness invariants and validation}
\label{sec:validation}

\subsection{The invariants}

The following invariants must hold after any change to the pipeline.

\begin{invariant}[Single Cartesian producer]
For a runtime chunk, every Cartesian vertex is written by
\code{OctreeTriangles.comp} stage 1.  No CPU code writes a chunk vertex.
The offline exporter is the only exception and is explicitly marked.
\end{invariant}

\begin{invariant}[Slot determinism]
The host assigns vertex slots in node-table order.  A rebuild with the same
octree revision produces the same slot for the same surface node.
\end{invariant}

\begin{invariant}[Span lifetime]
A draw entry references a span only after the finalize fence has signaled
and the drop counters are zero.  The old span is not freed until the cull
frame that could still reference it has retired.
\end{invariant}

\begin{invariant}[Index locality]
Every index written by stage 1 is a chunk-local slot
$s\in[0,\text{vertexCapacity})$ and is rebased by
\code{vertexOffset=baseVertex} at draw time.
\end{invariant}

\begin{invariant}[Owner consistency]
For a shared edge, both neighbouring chunks select the same owner, sample
the same four quadrants, and emit the same sub-segment list.  The owner
key is global (size, min, node index).
\end{invariant}

\begin{invariant}[Winding oracle]
Triangle winding is derived from the producing cell's corner signs, not
from a QEF normal.  A degenerate owner interpolation falls back to the
producing cell's signs.
\end{invariant}

\begin{invariant}[No silent truncation]
A break-list, stack, key-capacity, or span overflow sets a counter and
fails the chunk.  It never publishes a partial span.
\end{invariant}

\subsection{Validation layers}

Development builds enable Khronos validation, synchronization validation,
best-practices validation, and, when requested, GPU-assisted validation.
The engine treats every validation message as a bug.  The most relevant
checks for this pipeline are:
\begin{itemize}
\item descriptor range and offset alignment;
\item storage-buffer access outside the bound range;
\item synchronization hazards between compute writes and vertex/index
  reads;
\item queue submission threading;
\item image layout transitions in the offscreen vegetation/SDF targets.
\end{itemize}

\subsection{GPU-assisted validation}

GPU-assisted validation instruments shader resource access.  It is
especially useful here because a wrong \code{relVertexBase} or a stale span
may not crash: it reads another chunk's vertex and renders a stretched
triangle.  The GPU-assisted run catches the out-of-range descriptor access
at draw time.  The pass is slow and is used as a targeted check, not as the
default.

\subsection{Self-check against the CPU walker}

A diagnostic mode can rebuild a chunk with the CPU walker and compare the
triangle sets.  The comparison is quantization-tolerant because the GPU
and CPU may differ in the last bits of a QEF solve.  The check is:
\begin{itemize}
\item same triangle count;
\item same set of quantized triangle centroids;
\item same material assignment;
\item vertex deviation below a tolerance.
\end{itemize}
A large deviation indicates a flattening error (wrong cell, wrong LoD,
wrong owner), not a floating-point difference.

\subsection{Determinism tests}

A chunk is built twice from the same revision.  The vertex spans must be
byte-identical and the index stream must be a permutation of the first.
The triangle set is compared order-insensitively.  This catches:
\begin{itemize}
\item an unstable table order;
\item a slot assignment that depends on a hash map iteration;
\item an atomic append that changes the triangle set (it should not).
\end{itemize}

\subsection{Synchronization validation}

The topology pass is checked for:
\begin{itemize}
\item compute$\rightarrow$host visibility before the counter read;
\item compute$\rightarrow$vertex/index visibility before the draw;
\item no read of the pool while a finalize dispatch is still writing;
\item no free of a span still referenced by an in-flight cull.
\end{itemize}
The vegetation copy is checked for the same pattern: placement fence
$\rightarrow$ host read $\rightarrow$ copy submission $\rightarrow$
publish $\rightarrow$ consolidation copy.

\subsection{The queue-threading bug}

A release run produced
\begin{quote}\ttfamily\small\raggedright
vkQueueSubmit2(): THREADING ERROR : object of type VkQueue is\\
simultaneously used...
\end{quote}
The cause was \code{vkDeviceWaitIdle()} running without the submission
mutex while a worker submitted a topology or upload command.  The fix:
\begin{itemize}
\item \code{deviceWaitIdle()} takes the global submit mutex;
\item cleanup stops the CPU pools before the first idle wait.
\end{itemize}
After the fix, a full scene load and shutdown produced no threading or
validation messages.  This is an example of why the synchronization
section is part of the algorithm, not an afterthought.

\subsection{Regression checklist}

Any change to the mesher should run:
\begin{enumerate}
\item a debug build;
\item a load of the reference scene with all effects;
\item a validation scan of \code{logs/run.log};
\item a GPU-assisted run on a small scene;
\item a CPU-vs-GPU self-check on a chunk with a LoD boundary;
\item a rebuild/delete/rebuild cycle on the same chunk id.
\end{enumerate}
If any of these fails, the change is not ready.
